Il termine \emph{dunkle Materie}, ovvero materia oscura, venne coniato da Zwicky in un articolo del 1933 per giustificare una forte discrepanza rilevata tra la massa dinamica e la massa osservabile di un ammasso. 
Oggi, a quasi cento anni di distanza, gli scienziati sono convinti che essa sia la componente dominante della materia del nostro Universo, ma non ne hanno ancora compreso la natura microscopica. Sono stati avanzati diversi modelli teorici; il più accreditato è quello della materia oscura fredda (\emph{Cold Dark Matter}, CDM), che riesce a descrivere con successo l'evoluzione dell'Universo su scale cosmologiche, ma presenta tensioni su scala galattica e sub-galattica rispetto ai risultati osservativi. Possibili proposte alternative comprendono la \emph{Warm Dark Matter} (WDM) e la \emph{Self Interacting Dark Matter} (SIDM). 

In questo elaborato si mostra come lo \emph{strong gravitational lensing}, ovvero l'effetto di lente gravitazionale forte, può essere utilizzato per porre vincoli osservativi sui diversi modelli proposti, ponendosi fra le tecniche più promettenti per lo studio della materia oscura su piccola scala.

Data la complessità del problema affrontato, all'approccio basato sull'inferenza bayesiana gerarchica vengono affiancate simulazioni in ambiente controllato, con l'obiettivo di sviluppare un'intuizione fisica dei fenomeni coinvolti, pur senza riprodurne tutte le complessità osservative. Nella parte finale di questa tesi vengono quindi analizzati gli effetti della variazione dei parametri fisici e geometrici del sistema, sia in assenza sia in presenza di sotto-strutture, evidenziando le differenti firme gravitazionali associate a diversi profili di densità e discutendo al contempo i limiti del modello adottato.